\chapter{L'arbitro, e i moltiplicatori di Rao contro l'aggiunto}
\label{ch:referee}

Finora il programma ha verificato soltanto sé stesso. Questo capitolo
chiama in causa due giudici esterni: un codice di progetto di ugelli
scritto da altre persone in un altro linguaggio, i cui campi di moto
interi vengono sottoposti alle nostre equazioni, e la teoria del 1958
dell'ugello ottimale fatta con carta e matita, le cui previsioni
misuriamo su una forma che il nostro ottimizzatore non ha prodotto.
Superare entrambi significa che il motore risolve il problema giusto,
e non semplicemente un problema coerente con sé stesso.

\section{L'arbitro torna in linea}

\subsection{Che cosa occorre sapere prima}

Un ugello \emph{ideale}, o \emph{a lunghezza minima}, dà un flusso di
uscita perfettamente uniforme e puramente assiale nella minima
lunghezza in grado di farlo (Capitolo~\ref{ch:nozzle}). Un
\emph{contorno ottimizzato in spinta} (il TOC di Rao~\cite{rao})
risponde a una domanda diversa: per una \emph{data} lunghezza, quale
forma produce la massima spinta? È più corto e più pieno, il suo
flusso di uscita è deliberatamente non uniforme, e compra più spinta
per chilogrammo di ugello. Qui entrambe le forme raggiungono lo
stesso \emph{rapporto di espansione} $\varepsilon = A_e/A_t = 4$:
area di uscita quattro volte l'area di gola, il numero che dice
quanto l'ugello espande il gas. Ovunque, le tolleranze sono
\emph{bande derivate} calcolate dalla discretizzazione anziché
scelte, e ogni test porta con sé i \emph{rigettatori}: ingressi
corrotti che deve far fallire (Capitolo~\ref{ch:method}; simboli
nella Tabella~\ref{tab:notation}).

\subsection{La domanda}

Tutto quanto precede verifica il nostro codice contro sé stesso e
contro forme chiuse. Dov'è il programma \emph{indipendente} capace di
arbitrare interi campi di moto?

\subsection{Come funziona, per esteso}

GENO~\cite{geno} è quel programma: un progettista di ugelli in
Fortran indipendente --- base di codice separata, discretizzazione
separata, backend termodinamico separato --- e, per decisione
permanente, mai modificato da noi e mai differenziato. Come
\emph{generatore} produce contorni e campi di riferimento, in
particolare l'ugello ottimizzato in spinta di cui la
Sezione~\ref{sec:o33} ha bisogno: una forma genuinamente ottimale
nella cui produzione il nostro ottimizzatore non ha avuto parte. Come
\emph{arbitro}, i suoi campi vengono messi sotto processo contro le
\emph{nostre} equazioni. Il nostro solutore costruisce un flusso su
una maglia di linee d'onda intersecanti, imponendo in ogni punto di
maglia un insieme fissato di relazioni algebriche --- il suo
\emph{processo elementare} (Capitolo~\ref{ch:moc}); quelle relazioni
vengono ora valutate sui valori nodali di GENO, dove devono valere
entro la banda derivata. Il controllo è un insieme di \emph{piedi}
deliberatamente sbagliati --- i punti a monte, sulle due linee
d'onda, da cui un punto di maglia è costruito --- che \emph{non}
devono valere. L'accordo significa allora che i due codici risolvono
le stesse equazioni cella per cella, e non soltanto che le loro
uscite si somigliano.

Il vecchio binario di GENO era anteriore alla toolchain di questo
host, quindi è stato ricompilato dai sorgenti; il caso di riferimento
rigenerato è il TOC di tipo~2 a $\varepsilon = 4$ ($874\times2885$
nodi).

\subsection{Risultati di riferimento}

\begin{itemize}[itemsep=2pt]
\item \textbf{Oracolo fra codici: PASS.} 218/218 celle campionate del
  campo rigenerato soddisfano le nostre relazioni di processo
  elementare dentro la banda derivata, con una discriminazione
  $84\times$ fra piedi caratteristici corretti e deliberatamente
  sbagliati.
\item \textbf{Chiuso l'ultimo SKIP del Mattone~1.} Un test del
  Mattone~1 era ineseguibile finché l'arbitro non veniva compilato
  qui. Ora gira dal vivo, confrontando il nostro contorno ideale con
  quello di GENO stazione per stazione: 62/62 campioni in banda,
  errore massimo $7.6\times10^{-9}$~m.
\item \textbf{Il controllo positivo su $f_2$ è atterrato} (l'unico
  test dichiarato in sospeso della Sezione~\ref{sec:corner}). $f_2$
  combina velocità e direzione locali del flusso, e la teoria lo
  obbliga a essere costante lungo una linea di un ugello ottimale ---
  la \emph{caratteristica terminale}, definita più avanti. Sul
  contorno genuinamente ottimale di Rao è lì costante con una
  dispersione relativa di $4.0\times10^{-3}$ mentre lungo la parete
  varia di 0.60 --- una separazione $149\times$. Il caso di
  riferimento \emph{troncato} ha correttamente mostrato assenza di
  costanza: tagliato corto, un contorno ottimale non contiene più la
  propria superficie di controllo.
\end{itemize}

\begin{figure}[htbp]
\centering
\includegraphics[width=0.85\textwidth]{fig_contours.pdf}
\caption{I due protagonisti, disegnati in scala: distanza assiale $x$
[m] contro raggio $y$ [m], asse di simmetria a $y = 0$, gola a
$x = 0$, $y = y_t = 1$~m (quindi qui metri e raggi di gola
coincidono). Entrambe le pareti raggiungono il rapporto di espansione
$\varepsilon = 4$; si guardino le \emph{lunghezze}. L'ugello ideale
(a lunghezza minima) del Mattone~1 arriva a $L = 7.3409$~m, il
contorno ottimizzato in spinta di Rao prodotto da GENO si ferma a
5.5~m: più corto per costruzione (massima spinta per una \emph{data}
lunghezza). Una verifica incrociata non pianificata: GENO ha
rifiutato una richiesta di TOC più lungo citando la lunghezza ideale
$L = 7.3409$ --- esattamente quella del nostro gemello, da un codice
che il nostro non l'ha mai visto.}
\label{fig:contours}
\end{figure}

\section{I moltiplicatori di Rao contro il nostro aggiunto, su un
contorno realmente ottimale}
\label{sec:o33}

\subsection{La domanda}

Rao risolse il problema dell'ugello ottimale nel 1958 con la carta e
il calcolo delle variazioni. Questo programma lo risolve per
\emph{differenziazione automatica}: la macchina propaga le derivate
attraverso ogni operazione che il solutore esegue, dando la
sensibilità esatta della spinta a ogni parametro di forma --- e
eseguita all'indietro (un \emph{aggiunto}) una sola passata le
restituisce \emph{tutte} in una volta, a circa il costo di una
risoluzione di flusso aggiuntiva. La voce aperta più antica del
registro (la sua C-O33 in~\cite{theory}) chiede se le due risposte
concordino, termine per termine, su un contorno realmente ottimizzato
in spinta.

\subsection{Come funziona, per esteso}

Il terreno d'incontro è il contorno ottimale di GENO, marciato dal
\emph{nostro} motore a parete assegnata
(Capitolo~\ref{ch:wallinput}). Quando qualcosa viene massimizzato
sotto vincoli, ogni vincolo entra nella soluzione con una costante
che ne dà il prezzo --- un \emph{moltiplicatore di Lagrange}. La
soluzione di Rao~\cite{rao} ne porta due: $\lambda_2$, il prezzo di
tenere fissa la portata, e $\lambda_3$, il prezzo di tenere fissa la
lunghezza. Ciascuno è falsificabile, perché la teoria nomina una
combinazione di grandezze locali del flusso che deve eguagliarlo in
ogni punto di una curva --- un \emph{primo integrale}, una quantità
che le equazioni tengono lì costante. Quella curva è la
\emph{caratteristica terminale}: l'ultima linea d'onda $C^{-}$
dell'ugello, dall'asse al labbro, la superficie di controllo che il
calcolo delle variazioni seleziona perché spinta e portata vengono
valutate attraverso di essa. Lì:
\begin{equation}
f_2 = \frac{W\cos(\theta-\alpha)}{\cos\alpha} = -\lambda_2 ,
\qquad
2\pi y\,\rho W^{2}\sin^{2}\theta\,\tan\alpha = -\lambda_3 ,
\label{eq:firstintegrals}
\end{equation}
con $W$ il modulo della velocità [m/s] (il $q$ della
Tabella~\ref{tab:notation}), $\theta$ l'angolo del flusso, $\alpha$
l'angolo di Mach, $\rho$ la densità, $y$ il raggio --- entrambi
\emph{costanti lungo quella curva} se e solo se il contorno è
ottimale. Il test campiona entrambe le quantità sui nodi della
caratteristica terminale di GENO e richiede costanza entro bande
derivate nel nostro campo \emph{e} in quello di GENO, più l'accordo
fra i codici sugli stessi nodi.

Il lato aggiunto verifica la \emph{stazionarietà}, l'enunciato del
calcolo per lo stare su un massimo: in un massimo la derivata in ogni
direzione \emph{ammessa} si annulla, quindi piccole variazioni
ammesse non costano nulla al primo ordine. Le direzioni ammesse sono
deformazioni di parete che tengono fissi la lunghezza ed entrambi gli
estremi --- sei protuberanze gaussiane con estremi bloccati (centri
da $0.20L$ a $0.85L$). Una protuberanza che sposti il labbro sonda
una direzione che la teoria non protegge (un ugello più grande paga
sempre un poco di più) e deve dare un gradiente grande; le stesse sei
protuberanze su una parete deliberatamente perturbata sono il
controllo positivo. Ogni proiezione AD è verificata contro differenze
finite, cioè contro il ricalcolo della spinta con la forma spostata a
mano.

\subsection{Risultati di riferimento}

\textbf{9/9 PASS su due stadi}
(\carrier{validation/a1\_o33\_toc.py}). \emph{Integrali}: il nostro
campo sui nodi della caratteristica terminale di GENO mostra $f_2$
costante con dispersione relativa $4.045\times10^{-3}$ (quella di
GENO: $4.04\times10^{-3}$), e il livello coincide fra i codici a meno
di 0.235~m/s su 2634.7~m/s --- $9\times10^{-5}$ relativo
(Figura~\ref{fig:f2}); $\lambda_3$ coincide punto per punto (la sua
costanza è onestamente dichiarata irrisolvibile a questa precisione
di superficie: $\lambda_3 \sim \sin^{2}\theta$ è piccolo e ripido
vicino all'asse). \emph{Stazionarietà}: i sei gradienti interni
stanno $1000\times$ sotto la direzione del labbro e $40\times$ sotto
i gradienti sulla parete perturbata, AD${}={}$FD entro l'${\sim}1\%$
ovunque, e il contorno ottimale batte quello perturbato,
$J({\rm TOC}) > J({\rm perturbato})$
(Figura~\ref{fig:stationarity}). Un rigettatore è stato esso stesso
corretto dalla fisica lungo la strada:
Appendice~\ref{app:f2rejector}. Lo scarico formale del registro
(YAML e cablaggio della suite) è un atto da eseguire al momento del
commit sul repository WSL; il contenuto numerico esiste.


\begin{figure}[htbp]
\centering
\includegraphics[width=\textwidth]{fig_rao_construction.pdf}
\caption{Come viene realmente costruita la superficie di controllo
ottimale di Rao, disegnata sul campo realmente ottimizzato in spinta
(assi: posizione assiale e raggio in metri; le tenui linee grigie
sono la maglia caratteristica del codice di riferimento). (a)~Due
famiglie caratteristiche che si \emph{incrociano}: il bordo del
nucleo (arancione), una $C^{-}$ che corre dalla fine dell'espansione
iniziale in B giù fino all'asse, e la superficie di controllo (blu),
una $C^{+}$ che corre da D sull'asse su fino al labbro E. Il punto
rosso è il punto di lancio, dove la superficie di controllo lascia la
regione che la gola già determina. Tutto ciò che sta dentro il nucleo
ombreggiato è marcia ordinaria, calcolata una volta sola. (b)~La
famiglia a un parametro: ogni punto di lancio sul bordo del nucleo dà
un ugello candidato diverso. Linea continua dentro il nucleo --- lo
stesso flusso noto per ogni candidato; tratteggiata oltre, dove
ciascuno genera il proprio flusso e finisce al proprio labbro. La
ricerca fa scorrere il punto di lancio lungo quel bordo. (c)~Perché
questa singola curva possa essere tracciata senza conoscere il flusso
attorno a essa: una caratteristica ordinaria ha due incognite da
avanzare e solo la relazione di compatibilità per avanzarle, quindi
ha bisogno di una compagna dell'altra famiglia; la condizione di
ottimalità di Rao fornisce la seconda relazione mancante, e il
sistema si chiude.}
\label{fig:raoconstruction}
\end{figure}

\begin{figure}[htbp]
\centering
\includegraphics[width=\textwidth]{fig_f2.pdf}
\caption{Il primo integrale, misurato. (a)~Dove vivono le due curve:
$x$ [m] contro $y$ [m] sull'ugello ottimizzato in spinta di GENO ---
la parete (grigio) e il tratto campionato della caratteristica
terminale, la sua ultima linea d'onda (blu). (b)~$f_2$ [m/s] contro
ascissa curvilinea [m] lungo ciascuna curva: sale del 64\% lungo la
parete --- lì nessun invariante, e la teoria dice che non deve
esserci --- e resta piatto lungo la caratteristica terminale.
(c)~Ingrandimento su quella linea piatta, stesse unità: il nostro
campo (cerchi) si sovrappone a quello di GENO (croci) sugli stessi
nodi. Piattezza a $4\times10^{-3}$ relativi con accordo fra codici a
$9\times10^{-5}$: la costante $-\lambda_2$ di Rao esiste in un flusso
calcolato, e due codici senza rapporto fra loro ne trovano lo stesso
valore.}
\label{fig:f2}
\end{figure}

\begin{figure}[htbp]
\centering
\includegraphics[width=0.8\textwidth]{fig_stationarity.pdf}
\caption{Stazionarietà, in scala logaritmica. Asse orizzontale: la
deformazione di parete applicata --- sei protuberanze gaussiane,
etichettate dal loro centro come frazione $x/L$ [--] della lunghezza
di parete, più la deformazione che sposta il labbro. Asse verticale:
$|\mathrm{d}J/\mathrm{d}a|$ [N/m], variazione di spinta per metro di
ampiezza della protuberanza, logaritmico. Sul contorno ottimale (blu)
il gradiente sta $40\times$ sotto le stesse protuberanze su una
parete deliberatamente perturbata (grigio) e $1000\times$ sotto la
direzione del labbro (rosso): piatto nelle direzioni ammesse, ripido
nell'unica direzione che la teoria non protegge. È così che un ottimo
appare all'aggiunto.}
\label{fig:stationarity}
\end{figure}
